% Generated draft: requires mathematical and source review.
\documentclass[11pt]{article}
\usepackage[margin=25mm]{geometry}
\usepackage{fontspec}
\setmainfont{Latin Modern Roman}
\usepackage{amsmath,amssymb,mathtools,unicode-math}
\setmathfont{Latin Modern Math}
\usepackage{enumitem,needspace,xurl,hyperref}
\setlength{\parindent}{0pt}
\setlength{\parskip}{5pt}
\setlength{\emergencystretch}{3em}
\newcommand{\N}{\mathbb{N}}
\newcommand{\Z}{\mathbb{Z}}
\newcommand{\Q}{\mathbb{Q}}
\newcommand{\R}{\mathbb{R}}
\newcommand{\C}{\mathbb{C}}
\newcommand{\F}{\mathbb{F}}
\newcommand{\A}{\mathbb{A}}
\newcommand{\PP}{\mathbb{P}}
\newcommand{\E}{\mathbb{E}}
\newcommand{\eps}{\varepsilon}
\newcommand{\dd}{\,\mathrm{d}}
\begin{document}
\section*{UnsolvedMath ID: 30001852}
\textbf{Source problem code:} catalogue entry
\subsection{Definitions and mathematical statement}
The supplied recurrence has a natural real-variable formulation, stated below conditionally because the record does not specify the parameter domains or their admissible restrictions. Under the interpretation
\[
(C_1,C_2,C_3)\in\C^3,\qquad \alpha,\beta\in\R,
\]
define \(e(t)=\exp(2\pi i t)\). The Schwartz space \(\mathcal S(\R;\C)\) consists of all smooth functions \(f:\R\to\C\) such that
\[
\sup_{t\in\R}(1+|t|)^m|f^{(k)}(t)|<\infty
\qquad\text{for every }m,k\in\Z_{\geq0}.
\]
For a fixed admissible parameter tuple, the question is whether there exists \(f\in\mathcal S(\R;\C)\), \(f\not\equiv0\), satisfying
\[
f(t+1)=C_1f(t-1)+\bigl(C_2e(\alpha t)+C_3e(\beta t)\bigr)f(t)
\qquad\text{for every }t\in\R.
\]
The intended admissible parameter set and the intended quantification over that set are unresolved. In particular, the record does not specify nonvanishing conditions on the coefficients, distinctness conditions on the frequencies, or arithmetic restrictions. Thus this conditional formulation does not determine the full original question about four-point configurations on three lines.

The associated matrix recurrence can be written explicitly. For \(x\in\R\) and \(n\in\Z\), put
\[
v_n(x)=\begin{pmatrix}f(x+n)\\ f(x+n-1)\end{pmatrix},\qquad
A_n(x)=\begin{pmatrix}
C_2e(\alpha(x+n))+C_3e(\beta(x+n))&C_1\\
1&0
\end{pmatrix}.
\]
Then the scalar equation gives
\[
v_{n+1}(x)=A_n(x)v_n(x)\qquad(n\in\Z).
\]
A precise bilateral interpretation of a rapidly decaying matrix solution is a sequence \((v_n)_{n\in\Z}\) satisfying this recurrence and
\[
\sup_{n\in\Z}(1+|n|)^N\|v_n\|<\infty
\qquad\text{for every }N\in\Z_{\geq0}.
\]
The record's second question can accordingly be read as asking whether there is a Lebesgue-measurable set \(E\subseteq\R\) of positive measure such that, for every \(x\in E\), the matrix recurrence admits a nonzero rapidly decaying bilateral solution. Here \(x\) is the base point, rather than the initial vector. A nonzero Schwartz solution supplies such solutions on a positive-measure set. The supplied claim that the two questions are equivalent requires further hypotheses or an argument: separate rapidly decaying orbit solutions do not by themselves specify a globally smooth Schwartz function. The required dependence on \(x\), compatibility between base points, and decay conventions are not supplied. No backward invertibility is assumed; the matrices are invertible exactly when \(C_1\ne0\).
\subsection{Short English statement}
The problem asks whether a nonzero Schwartz function can satisfy the displayed translation--modulation recurrence, motivated by four-point configurations on three lines. The parameter restrictions are missing, and the claimed equivalence with rapidly decaying matrix solutions on a positive-measure set needs clarification. This is a conditional statement-only draft requiring human mathematical review.
\subsection{Sources}
\begin{itemize}
\item Local UnsolvedMath registry: supplied complete record for permanent ID \(30001852\), source problem code \texttt{catalogue entry}; project registry \texttt{data/unsolvedmath/problems.json}, with \texttt{catalog/problems.json} designated as the supplementary registry. The original wording is: ``How can four-point configurations lying on three rather than two lines be handled?'' No individual proposer is named; the supplied proposed year is 2011.
\item Supplied report citation: \emph{Real Analysis, Harmonic Analysis and Applications} (2012), Oberwolfach Reports, pp.\ 2045--2111, DOI \texttt{10.4171/owr/2011/36}. \url{https://doi.org/10.4171/owr/2011/36}. The record attributes the problem to workshop open-problem discussions. Its verification field says that the official report was checked, the concrete recurrence for the four-point, three-line HRT configuration was restored, and the exponential convention follows the report's preceding time-frequency formula. These are supplied verification claims, not independent verification performed for this draft.
\item Supplied literature source, labeled ``Oussa four point counterexample'': \url{https://arxiv.org/abs/2608.07604}. No further bibliographic details are supplied.
\item Supplied literature source, labeled ``Oussa certified four point results'': \url{https://arxiv.org/abs/2604.21228}. No further bibliographic details are supplied.
\item Supplied literature source, labeled ``OWR report 36'': \url{https://doi.org/10.4171/owr/2011/36}.
\item Literature assessment dated exactly \texttt{2026-08-21}: the supplied status is ``partially solved,'' with the assessment that the exact three-line four-point case is not classified, that 2026 work settles some four-point cells, and that another cell has counterexamples. This is supplied dated curation derived from the OWR curation, not an independently verified mathematical or literature assessment. The record expressly requests independent verification of time-sensitive status and cited claims.
\end{itemize}
\end{document}
